\documentclass[twocolumn]{article}

\usepackage{hyperref} % references to other parts of the paper if clicked

\usepackage{listings} % putting code in paper
\renewcommand\lstlistingname{Source Code}
\renewcommand\lstlistlistingname{Source Code}

% code color style
\usepackage{xcolor}
\definecolor{dkgreen}{rgb}{0,0.6,0}
\definecolor{dred}{rgb}{0.545,0,0}
\definecolor{dblue}{rgb}{0,0,0.545}
\definecolor{lgrey}{rgb}{0.9,0.9,0.9}
\definecolor{gray}{rgb}{0.4,0.4,0.4}
\definecolor{darkblue}{rgb}{0.0,0.0,0.6}
\lstset{ 
    backgroundcolor=\color{lgrey},  
    basicstyle=\footnotesize \ttfamily \color{black} \bfseries,   
    breakatwhitespace=false,       
    breaklines=true,               
    captionpos=b,                   
    commentstyle=\color{dkgreen},   
    deletekeywords={...},          
    escapeinside={\%*}{*)},                  
    frame=single,                  
    keywordstyle=\color{purple},  
    morekeywords={BRIEFDescriptorConfig,string,TiXmlNode,DetectorDescriptorConfigContainer,istringstream,cerr,exit}, 
    identifierstyle=\color{black},
    stringstyle=\color{blue},      
    language=C++,                
    numbers=right,                 
    numbersep=5pt,                  
    numberstyle=\tiny\color{black}, 
    rulecolor=\color{black},        
    showspaces=false,               
    showstringspaces=false,        
    showtabs=false,                
    stepnumber=1,                   
    tabsize=5,                     
    title=\lstname,                 
}

% for piece-wise functions
\usepackage{amsmath}

% decreasing the left and right margins
\usepackage{blindtext}
\usepackage{geometry}
\geometry{
    a4paper,
    total={170mm,260mm},
    left=15mm,
    right=15mm,
    top=20mm,
}

\usepackage{graphicx} % can add images to paper
\usepackage{float} 
\usepackage{adjustbox} % adjusting pictures
\graphicspath{ {pictures/} }

\begin{document}

% ============= Mutex =============
\section{Mutex (\underline{Mut}ual \underline{Ex}clusion)}

\underline{Definition}: class is a synchronization primitive that can be used to protect shared data from being simultaneously accessed by multiple threads.

\bigbreak \noindent \underline{Libraries Required}: $<$mutex$>$

\subsection{Lock}

\subsection{Try Lock}

\subsection{Unlock}

% ============= Lock Guard(s) =============
\section{Lock Guard(s)}
\underline{Definition}: mutex wrapper that provides a convenient RAII-style mechanism for owning a mutex for the duration of a scoped block.

\bigbreak \noindent \underline{Libraries Required}: $<$mutex$>$

\noindent Different types of lock guards:
\begin{enumerate}
    \item lock guard
    \item scoped lock
    \item unique lock
    \item shared lock
\end{enumerate}

% ============= Thread =============
\section{Thread} \label{sec:thread}
\underline{Definition}: A lightweight sub-task or flow of control that lives inside a process.

\subsection{Join}

\subsection{Detach}

% ============= JThread =============
\section{JThread}
\underline{Definition}: a modern replacement for std::thread that represents a single thread of execution. The "j" stands for joining, which directly highlights its core purpose: it implements the RAII (Resource Acquisition Is Initialization) pattern for managing thread lifecycles.

% ============= Process =============
\section{Process}
\underline{Definition}: An active instance of a program with its own isolated memory space and system resources 

% ============= Async =============
\section{Async}
\underline{Definition}: a mechanism the function f asynchronously (potentially in a separate thread which may be part of a thread pool) and returns a std::future that will eventually hold the result of that function call.

\subsection{Launch Policy}
Launch policies allow the programmer when to trigger the async operation(s).

\subsubsection{launch::async}
Launches a new thread to call your given \textbf{function} as if a thread object 

\subsubsection{launch::deferred}
Your given \textbf{function} is deferred until the shared state of the returned future is accessed with \underline{wait} or \underline{get} (member functions in a future).

% ============= Future =============
\section{Future}
\underline{Definition}: is a class template that acts as a placeholder or a handle for a value (or exception) that will eventually be computed by an asynchronous operation.

\bigbreak \noindent \underline{Libraries Required}: $<$future$>$

% ============= Promise =============
\section{Promise}
\underline{Definition}: is a mechanism to store a value or an exception that will be asynchronously retrieved by another thread.

\bigbreak \noindent \underline{Libraries Required}: $<$future$>$

% ============= Package Task =============
\section{Package Task}
\underline{Definition}: class template that wraps any callable target (such as a function, lambda expression, or bind expression) so that its return value or thrown exception can be retrieved asynchronously via a std::future.

\bigbreak \noindent \underline{Libraries Required}: $<$future$>$


% ============= Atomics =============
\section{Atomics}
\underline{Definition}: operations or types that guarantee indivisible execution and inter-thread synchronization without the use of explicit locks (like mutexes).

\bigbreak \noindent \underline{Libraries Required}: $<$atomic$>$

% ============= Stop Token =============
\section{Stop Token}
\underline{Definition}: is a thread-safe, lightweight C++20 class that provides a read-only view of a shared stop-state to check whether an asynchronous operation or a std::jthread has been asked to stop.

\bigbreak \noindent \underline{Libraries Required}: $<$stop\_token$>$

% ============= Condition Variables =============
\section{Condition Variables}
\underline{Definition}: a mechanism that allows one or more threads to be suspended (put to sleep) until another thread modifies a shared state and notifies them.

\bigbreak \noindent \underline{Libraries Required}: $<$condition\_variable$>$

\end{document}